\documentclass[12pt]{beamer}
%\documentclass[12pt,handout]{beamer}



\usepackage{hyperref}
\usepackage{minted}
\usepackage[round]{natbib}

\hypersetup{
    colorlinks=true,
    linkcolor=blue,
    filecolor=blue,      
    urlcolor=blue,
    citecolor=blue,
}

% If you want to print handouts
%\documentclass[handout]{beamer}

% Changing the fonts: this will make the slides more readable and the math look like regular tex math
%\usefonttheme{arial}

% Titles will appear in Small Cap Serif
\usefonttheme{structuresmallcapsserif}

% -----------------------------------------
% Center the Frame Title
% -----------------------------------------

\setbeamertemplate{frametitle} {
\begin{centering}
\vspace{0.1in} \insertframetitle
\par
\end{centering}}


% -----------------------------------------
% Number the slides
% -----------------------------------------

\setbeamertemplate{footline}[frame number]

% -----------------------------------------
% Get rid of the irritating navigation bar
% -----------------------------------------

\setbeamertemplate{navigation symbols}{}
% -----------------------------------------
% Begin Document
% -----------------------------------------
\title{\Large\textsc{Media and government censorship: Evidence from Mexico}}
\author[Zago] % Esto es lo que sale en el footer
{\textbf{Eduardo Zago}\inst{1}  \vspace{10pt}}

\institute []% Esto es lo que sale en el footer
{  \inst{1}   PhD Student NYU Politics}

\date{\textsc{}}


\begin{document}
\maketitle
%  I separate each slide with a line:

% -----------------------------------------

%  I start with the screen black, like using the b key in Powerpoint



% -----------------------------------------
%  Restore normal white background

\beamersetaveragebackground{white}


% -----------------------------------------

%  Example of a slide with animated math:
\section{Motivation}

\begin{frame}
\frametitle{Context and Research Question}

\small
\begin{itemize}
\item Media outlets have been identified as a critical component of electoral accountability in Latin America \citep{Audits-Ferraz-QJE, PolInfo-Enriquez-JEEA, Marshall}.
\item Even in the absence of outright censorship, the government still has techniques to silence critics \citep{Censor-Pratt-AER, Autocracy-Guriev-JPE}
\item In Mexico, there is anecdotal evidence that past administrations use the social communication spending for this purposes \citep{Times}.

\textbf{RQ1: How does the government react to media outlets being critic of it?}

\textbf{RQ2: How does the media react to censorship or critics coming from the President?}



\end{itemize}
\end{frame}

\begin{frame}
\frametitle{Research Design RQ1}

\small
\begin{itemize}
    \item Estimate the impact of critical coverage on government communication spending by media outlet.
    \item OLS probably biased due to time-based endogeneity and unobserved confounding factors related to criticism level and resource allocation.
    \item \textbf{Instrumental Variable (IV) Approach:} Use lottery assignment to the first, second or third row at morning press conferences as an instrument for $NumCr_{i, t}$.
\item First Stage Equation:

$$NumCr_{i, t} = \theta_0 + \theta_1 FirstRow_{i, t} + \mathbf{X'}_{i,t}\gamma + \epsilon_{i, t}$$
\end{itemize}
\end{frame}

\begin{frame}
\frametitle{Feedback}

\small
\begin{itemize}
    \item Any possible threats to the exclusion restriction that I should think about?
    \item 
    \item I am struggling thinking about the period of analysis. 
    \item If month then the instrument should be the number of times the media outlet was assigned to the first row that month. 
    \item This probably affects the critical articles of that month, but the government spending? Or in T+1.

    \item Any ideas on the scrapping of newspapers would be greatly appreciated.

\end{itemize}
\end{frame}




\begin{frame}{References}
\scriptsize
    \bibliographystyle{apsr}
    \typeout{}

    \singlespacing

    \bibliographystyle{apsr}
    \bibliography{literature}
\end{frame}

\section{Appendix}

\begin{frame}
\frametitle{Research Design RQ2}

\small
\begin{itemize}
\item Narrative comparison between what the President said everyday in the press conference and what the media covered.
\item H1: media reacts to attacks from the President by cherry-picking on the conference and switching the narrative.

\item Use \cite{QNarr-Waight-WP} to extract all narratives from both sources, observationally compare them

\item Heterogeneous analysis by media outlet ideology, topic discussed (violence related, controversial topics), etc.

\end{itemize}
\end{frame}

\end{document}